\section{Burst size and burst time}
\label{bdc:sec:burst}

The moments just derived already tell us something about the rupture event
itself. There are two quantities of interest: the burst size
$\Kb=X_{\tau-}$, which is the load carried out of the cell, and the burst time
$\tau$. Neither is a marginal moment of $X_t$ alone, but both can be reached
from the functions now in hand.

\subsection{Conditioning on rupture}
\label{bdc:sec:conditioning}

The mean release $V(t)$ includes paths which have died internally and will never
release anything. To obtain the mean load released by cells which have ruptured
by time $t$, split the expectation according to the event $\{\tau\le t\}$:
\[
\ex{W_t}
=\ex{W_t\mid\tau>t}\,I(t)
+\ex{W_t\mid\tau\le t}\,[1-I(t)].
\]
The first conditional expectation is zero because $W_t=0$ before rupture.
Hence
\[
\ex{W_t\mid\tau\le t}=\frac{V(t)}{1-I(t)}.
\]
Letting $t\to\infty$ conditions on eventual rupture, irrespective of when it
occurs.

\begin{corollary}[The mean burst size]
\label{bdc:cor:burstmean}
Conditional on eventual rupture,
\begin{equation}
\label{bdc:eq:burstmean}
\ex{\Kb\mid\text{burst}}
=\lim_{t\to\infty}\frac{V(t)}{1-I(t)}
=\frac{V_\infty}{1-b}
=\frac{\lambda-\mu}{\delta}+\frac{1}{1-b}
=\frac{a}{a-1}.
\end{equation}
When $\mu=0$, every realisation eventually ruptures and this conditional mean
equals $V_\infty$.
\end{corollary}

The second moment follows in the same way. As $t\to\infty$, we have $I\to b$
and $K\to0$, so \cref{bdc:eq:Wsq} gives
\[
\ex{W_\infty^2}
=\frac{2(\lambda-\mu)}{\delta}V_\infty
+\frac{\lambda+\mu}{\delta}(1-b)+1.
\]
Dividing by the rupture probability $1-b$, subtracting the square of
\cref{bdc:eq:burstmean}, and using \cref{bdc:eq:identities} yields the variance.

\begin{corollary}[The first two moments of the burst law]
\label{bdc:cor:burstvar}
Conditional on rupture, the burst size satisfies
\begin{equation}
\label{bdc:eq:burstvar}
\ex{\Kb\mid\text{burst}}=\frac{a}{a-1},
\qquad
\var{\Kb\mid\text{burst}}=\frac{a}{(a-1)^2},
\end{equation}
for every $\mu\ge0$.
\end{corollary}

These coincide with the first two moments of a geometric distribution on
$\{1,2,\ldots\}$ with ratio $1/a$, or success probability $(a-1)/a$.
The agreement points towards the geometric law, but two moments do not identify
a distribution. \fwd{singlecell}{Chapter~5} derives the full burst-size law and
shows that the indication is correct.

The no-death case gives a convenient check. With $\mu=0$, we have
$b=0$, $I(\infty)=0$, $K(\infty)=0$ and
$V_\infty=1+\lambda/\delta$. Equation~\cref{bdc:eq:Wsq} then gives
\[
\ex{W_t^2}\Big|_{t\to\infty}
=\frac{2\lambda}{\delta}
\Bigl(1+\frac{\lambda}{\delta}\Bigr)
+\frac{\lambda}{\delta}+1
=\frac{2\lambda^2}{\delta^2}+\frac{3\lambda}{\delta}+1.
\]
This is $(2-s)/s^2$ for
$s=\delta/(\lambda+\delta)=(a-1)/a$, the second moment of that geometric law.

\begin{figure}[!htbp]
\centering
\includegraphics[width=0.72\textwidth]{F4_burst_size_preview.pdf}
\caption{The burst-size law in the no-death case:
$\pr{\Kb=k}=(\delta/\lambda)a^{-k}$, a geometric law on
$\{1,2,\ldots\}$ with ratio $1/a$. Here
$\lambda=1$, $\mu=0$, $\delta=0.1$ and $a=1.1$.}
\label{bdc:fig:burstsize}
\end{figure}

\Cref{bdc:fig:burstsize} previews the distribution in the case $\mu=0$, when
every path bursts and no conditioning mass is lost to internal extinction.

\subsection{The burst time}
\label{bdc:sec:bursttime}

The distribution of $\tau$ is contained in $I$:
\[
\pr{\tau\le t}=\pr{X_t=H}=1-I(t).
\]
Its density on finite times is therefore
\begin{equation}
\label{bdc:eq:bursttime}
\varphi(t)=-\ddt{I}=\delta J(t).
\end{equation}
The identity has a simple reading. Since any unit can trigger rupture, the
instantaneous rupture density is proportional to the current mean load.

When $\mu>0$, this density is defective:
\[
\int_0^\infty\varphi(t)\,\dd t
=1-\lim_{t\to\infty}I(t)
=1-b<1.
\]
The missing mass $b$ belongs to cells which die internally and consequently
have $\tau=\infty$. Only in the no-death case is $\varphi$ a proper density on
$[0,\infty)$. The full law of $\tau$, and its joint law with $\Kb$, are derived
in \fwd{singlecell}{Chapter~5}.

\subsection{The conditional mean and its plateau}
\label{bdc:sec:plateau}

There is a companion to the mean burst-size calculation. Although the
unconditional load $J(t)$ tends to zero, the mean among cells which are still
productive,
\[
\frac{J(t)}{\Ifix(t)},
\]
approaches the finite value $a/(a-1)$. This is the same value as the conditional
mean burst size in \cref{bdc:eq:burstmean}. The coincidence is suggestive:
late-rupturing cells release the load they carry, and the conditional load itself
is approaching a plateau. It is not, by itself, proof that the two limiting
distributions agree. That stronger statement belongs to the quasi-stationary
analysis \cite{yaglom1947certain,Meleard2012}; \fwd{singlecell}{Chapter~5}
establishes it by deriving the limiting conditional law.

\Cref{bdc:fig:condmean} separates the no-death and death cases. Without death,
conditioning on non-catastrophe is the same as conditioning on productivity.
With death, an intact cell may already be empty, and the two normalisations no
longer agree.

\begin{figure}[!htbp]
\centering
\includegraphics[width=0.94\textwidth]{F4_conditional_mean.pdf}
\caption{Conditional mean loads. (a) For $\mu=0$, $\lambda=1$ and
$\delta=0.1$, $J(t)/I(t)$ and $V(t)$ tend to
$1+\lambda/\delta=11$. (b) For $\lambda=1$, $\mu=0.2$ and $\delta=0.05$,
$J(t)/\Ifix(t)$ approaches $a/(a-1)$, whereas $J(t)/I(t)$ tends to zero
because non-catastrophe also includes cells that died internally.}
\label{bdc:fig:condmean}
\end{figure}

\FloatBarrier
